\clearpage
\section{Seven short exercises}\label{c:exercises}
The first five exercises stay with \(G=\Z/6\Z\),
\(A=\{0,3\}\), \(B=\{0,2,4\}\). They test the maps, not just the final
factor list. Worked answers follow on the next two pages.

\begin{enumerate}
\item \textbf{Direction and witnesses.}
For \(R=r_2s_1\), give an intermediate witness showing \(2R5\).
Does \(5R2\)? Describe \(R^\op=s_1r_2\).
Explain why associativity of relation composition does not allow us
to interchange \(r_2\) and \(s_1\).

\item \textbf{The two ambiguities.}
Compute \(R^\op R\) and \(RR^\op\), including their supports.
Why is the first not an equivalence relation on all of \(G\)?
Write the induced quotient isomorphism and its inverse.

\item \textbf{Three corners.}
Let \(T\leq X\times Y\) be a subgroup, with \(X,Y\) possibly
noncommutative. Verify the ordered product that completes
\((x,y),(x',y),(x',y')\) to \((x,y')\).
For the congruence modulo \(B\), apply it to
\((0,2),(4,2),(4,4)\).
Give an arbitrary subset of \(C_2\times C_2\) for which the
three-corner property fails.

\item \textbf{What saturation repairs.}
Let \(T=\{(b,b):b\in B\}\subseteq B\times G\), and let
\(p=\id_B\times q_A\).
Show \(p(T)=p(R)\), but \(p^{-1}(p(T))=R\ne T\).
Compute \(r_2T\) and \(r_2R\). Which relation includes every representative
of each output coset?

\item \textbf{Reconstruct the rectangle.}
Without looking at §\ref{c:rectangle}, compute all four output quotients
and all four input quotients. Recover the two refined chains.
Verify that \(x+B\mapsto3x\) is inverse to \(a\mapsto a+B\).
Explain the different supports of the two trivial cells.

\item \textbf{Normal in the next term.}
Let \(A_4\) be the group of even permutations of four symbols, composed
from right to left. Inside \(A_4\), let
\[
 H=\{1,(12)(34)\},\qquad
 V_4=\{1,(12)(34),(13)(24),(14)(23)\}.
\]
Show that \(1<H<V_4<A_4\) is a composition series, but
\(H\) is not normal in \(A_4\).
For the last assertion, conjugate \((12)(34)\) by \((123)\).
Why would a proof considering only ambient-normal subgroups miss this
series?

\item \textbf{What the factors forget.}
Give composition series for \(C_4\) and \(C_2\times C_2\) with the same
simple factors. Prove the groups are not isomorphic.
Then count the homomorphisms \(C_2\to C_4\) and
\(C_2\to C_2\times C_2\). What information does this small probe detect
that the factor list does not?
\end{enumerate}

\route{Exercises 1--5 check the comparison mechanism. Exercises 6--7
expose two features deliberately invisible in the easy \(C_6\) example:
non-ambient normality and different ways to assemble the same factors.}

\clearpage
\section{Worked answers}\label{c:answers}
\textbf{1. Direction and witnesses.}
The intermediate point is \(2\): \(2s_12\) and \(2r_25\), since
\(5-2=3\in A\). But \(5\) is outside the input support \(B\), so \(5R2\)
is false. The opposite relation has pairs
\[
 R^\op=\{(b,b),(b+3,b):b\in B\}.
\]
It has input support \(G\) and output support \(B\). First move by \(A\),
then require the endpoint to lie in \(B\).
Associativity permits regrouping three successive operations; it does
not reverse their order. Here \(r_2s_1\ne s_1r_2\), as the supports
already show.

\medskip
\textbf{2. The two ambiguities.}
Distinct rows in the relation table have disjoint outputs, so
\[
 R^\op R=\Delta_B.
\]
It is reflexive on \(B\), but not on \(G\): for example \((1,1)\) is
missing. The other composite is
\[
 RR^\op=
 \{0,3\}^2\ \cup\ \{2,5\}^2\ \cup\ \{4,1\}^2=r_2.
\]
Its support is \(G\), and its three classes are the \(A\)-cosets.
The quotient map is \(B\to G/A,\ b\mapsto b+A\).
Its inverse \(x+A\mapsto4x\) is well defined because \(4A=0\).
The two identities \(4b=b\) on \(B\) and \(4x-x\in A\) verify both
composites directly.

\medskip
\textbf{3. Three corners.}
In the direct product group,
\[
\begin{split}
 (x,y)(x',y)^{-1}(x',y')
 &=\bigl(xx'^{-1}x',\,yy^{-1}y'\bigr)\\
 &=(x,y').
\end{split}
\]
The subgroup contains this product. Nothing is commuted, so the argument
works for arbitrary groups. In the additive \(C_6\) example the same
operation reads
\[
 (0,2)-(4,2)+(4,4)=(0,4),
\]
which indeed belongs to the congruence modulo \(B\).

The subset
\(\{(0,0),(1,0),(1,1)\}\subseteq C_2\times C_2\)
contains three corners but omits \((0,1)\).
It is not a subgroup. Merely having a subset of a product supplies no
three-corner property.

\pausehere{The ternary expression \(uv^{-1}w\) need not be a group
homomorphism in its three variables. The proof uses closure of the
subgroup under the displayed operations, not such a homomorphism claim.}

\clearpage
\textbf{4. What saturation repairs.}
Both images are the graph
\[
 L=\{(0,A),(2,2+A),(4,4+A)\}\subseteq B\times G/A.
\]
A pair \((b,y)\) lies in \(p^{-1}(L)\) exactly when \(y+A=b+A\),
or \(y-b\in A\). This is precisely \(R\). Therefore
\(r_2T=R\) and \(r_2R=R\). Only \(R\) includes all representatives
of each output coset; \(T\) includes just one.

\medskip
\textbf{5. Reconstruct the rectangle.}
The output and input quotient tables, respectively, are
\[
 \begin{pmatrix}0/0&A/0\\G/A&G/G\end{pmatrix},
 \qquad
 \begin{pmatrix}0/0&G/B\\B/0&G/G\end{pmatrix}.
\]
Their row and column chains are
\[
 0\subseteq0\subseteq A\subseteq G\subseteq G,\qquad
 0\subseteq0\subseteq B\subseteq G\subseteq G.
\]
The formula \(x+B\mapsto3x\) is well defined since \(3B=0\).
For \(a\in A\), \(3a=a\); for \(x\in G\), \(3x-x=2x\in B\).
These prove both inverse identities. The first trivial cell has support
\(\{0\}\), and the last has support \(G\). Each quotient has one element,
although their supports and relations differ.

\medskip
\textbf{6. Normal in the next term.}
\(V_4\) is abelian, so \(H\nml V_4\). Conjugation by an element of
\(A_4\) permutes the three double transpositions, so \(V_4\nml A_4\).
The successive factors have orders \(2,2,3\), and every group of prime
order is cyclic and simple. These steps give a composition series.
However,
\[
 (123)(12)(34)(123)^{-1}=(23)(14)\notin H.
\]
Thus \(H\) is not normal in \(A_4\). The theorem concerns subnormal
series, where each term is normal in the next; it cannot be reduced to
chains in the lattice of subgroups normal in the whole ambient group.

\medskip
\textbf{7. What the factors forget.}
Use
\[
 0<\{0,2\}<\Z/4\Z,\qquad
 0<C_2\times0<C_2\times C_2.
\]
Both factor lists are \(C_2,C_2\). The first group has an element of order
four, while the second does not, so they are not isomorphic.
A homomorphism out of \(C_2\) is determined by the image of its generator,
which must be killed by two. There are two possibilities in \(C_4\)
and four in \(C_2\times C_2\). The probe detects their different
two-torsion subgroups. The simple factors alone do not recover how those
factors are assembled.
